\documentclass[../DoAn.tex]{subfiles}
\begin{document}

Chương này trình bày cơ sở lý thuyết và công nghệ được sử dụng để xây dựng hệ thống \textit{multi-tenant RAG chatbot hỗ trợ tư vấn bán hàng, so sánh và tham chiếu giá sản phẩm nội thất}. Nội dung chương được chia thành hai nhóm: cơ sở lý thuyết phục vụ bài toán và các công nghệ dùng để triển khai hệ thống.

\section{Cơ sở lý thuyết}
\label{section:3.1}

\subsection{Mô hình multi-tenant}
\label{subsection:3.1.1}

Hệ thống phục vụ nhiều cửa hàng nội thất trên cùng một nền tảng, trong đó mỗi cửa hàng có dữ liệu sản phẩm, chatbot, knowledge base, hội thoại, lead, yêu cầu mua hàng và thành viên riêng. Vì vậy, hệ thống sử dụng mô hình multi-tenant\footnote{Microsoft Learn, \textit{Tenancy models for a multitenant solution}, \url{https://learn.microsoft.com/en-us/azure/architecture/guide/multitenant/considerations/tenancy-models}, truy cập lần cuối ngày 29/05/2026} để nhiều cửa hàng dùng chung hệ thống nhưng vẫn tách biệt dữ liệu và quyền truy cập.

Trong phạm vi đồ án, mỗi tenant tương ứng với một cửa hàng nội thất. Multi-tenant là mô hình một ứng dụng phục vụ nhiều tenant khác nhau, đồng thời cần tách biệt tài nguyên, dữ liệu và cấu hình giữa các tenant.

Trong hệ thống, các dữ liệu nghiệp vụ quan trọng được gắn với \texttt{tenant\_id}. Backend sử dụng thông tin tenant để giới hạn dữ liệu khi quản lý chatbot, knowledge base, hội thoại, lead và yêu cầu mua hàng, giúp người dùng của cửa hàng này không truy cập được dữ liệu của cửa hàng khác.

\subsection{Retrieval-Augmented Generation và knowledge base}
\label{subsection:3.1.2}

Chatbot tư vấn bán hàng cần trả lời dựa trên dữ liệu sản phẩm của cửa hàng, vì vậy hệ thống sử dụng Retrieval-Augmented Generation (RAG)\footnote{Lewis và cộng sự, \textit{Retrieval-Augmented Generation for Knowledge-Intensive NLP Tasks}, \url{https://arxiv.org/abs/2005.11401}, truy cập lần cuối ngày 29/05/2026}. RAG kết hợp truy xuất thông tin từ nguồn tri thức bên ngoài và đưa dữ liệu liên quan vào ngữ cảnh để mô hình ngôn ngữ tạo câu trả lời.

Trong đồ án, knowledge base là tập dữ liệu phục vụ chatbot, gồm thông tin sản phẩm, mô tả, giá, loại sản phẩm, chất liệu, kích thước, phong cách và các tài liệu liên quan do cửa hàng cung cấp. Knowledge base được tổ chức theo tenant, nên chatbot của cửa hàng nào chỉ truy xuất dữ liệu của cửa hàng đó.

Quy trình RAG gồm: nhận câu hỏi người dùng, xác định tenant và ngữ cảnh hội thoại, truy xuất dữ liệu liên quan từ knowledge base, gửi dữ liệu sang AI Service và tạo câu trả lời bằng Claude Sonnet.

\subsection{Dữ liệu sản phẩm, so sánh sản phẩm và tham chiếu giá}
\label{subsection:3.1.3}

Dữ liệu sản phẩm trong hệ thống gồm các thuộc tính chính như tên, loại sản phẩm, giá, chất liệu, kích thước, phong cách, mô tả và cửa hàng cung cấp. Các thông tin này được dùng để chatbot tư vấn, so sánh sản phẩm và tham chiếu giá.

Với chức năng so sánh sản phẩm, hệ thống tổng hợp các sản phẩm được người dùng nhắc đến và đối chiếu theo các tiêu chí như giá, chất liệu, kích thước, loại sản phẩm, phong cách và mức độ phù hợp với nhu cầu. Nếu thiếu dữ liệu ở tiêu chí nào, hệ thống cần nêu rõ phần còn thiếu.

Với chức năng tham chiếu giá, hệ thống chỉ đưa ra nhận xét trong phạm vi dữ liệu hiện có, không nhằm định giá toàn bộ thị trường. Ví dụ, giá của một mẫu bàn làm việc có thể được so sánh với các sản phẩm cùng loại trong dữ liệu của hệ thống.

\subsection{Nhận diện khách hàng liên kênh}
\label{subsection:3.1.4}

Trong hệ thống chatbot bán hàng, một khách hàng có thể tương tác qua nhiều kênh như Messenger hoặc Telegram. Mỗi kênh có định danh ngoài riêng, ví dụ Messenger dùng \texttt{PSID}, Telegram dùng \texttt{chat\_id} hoặc \texttt{user\_id}. Nếu chỉ lưu theo từng kênh, cùng một khách hàng có thể bị tách thành nhiều hồ sơ.

Trong phạm vi đồ án, hệ thống liên kết các định danh ngoài về một hồ sơ khách hàng nội bộ khi có đủ căn cứ. Việc gộp khách hàng không dựa trên tên hiển thị vì thông tin này dễ trùng hoặc thay đổi, mà ưu tiên số điện thoại hoặc email đã được chuẩn hóa.

Cách này giúp cửa hàng nhận biết cùng một khách hàng trên nhiều nền tảng và hạn chế gộp nhầm hồ sơ.

\subsection{Kiểm soát prompt, context và token trong RAG}
\label{subsection:3.1.5}

Trong hệ thống RAG, không đưa toàn bộ dữ liệu sản phẩm hoặc knowledge base vào prompt. Hệ thống chỉ truy xuất các chunk hoặc sản phẩm liên quan đến câu hỏi hiện tại, đồng thời giới hạn số lượng và độ dài context để tránh vượt token, tăng độ trễ và nhiễu thông tin.

Trước khi gọi mô hình ngôn ngữ, hệ thống xác định đúng tenant, chatbot, chế độ hội thoại và phiên bản knowledge base. Với dữ liệu sản phẩm, context chỉ giữ các trường cần thiết như tên, giá, chất liệu, kích thước, mô tả ngắn và URL nguồn.

Prompt được tổ chức theo template gồm chỉ dẫn hệ thống, context truy xuất, lịch sử hội thoại cần thiết và câu hỏi hiện tại. Cách này giúp mô hình trả lời đúng dữ liệu được cung cấp, không tự tạo thông tin ngoài context và nêu rõ khi thiếu dữ liệu.


\section{Công nghệ sử dụng}
\label{section:3.2}

\subsection{Claude Sonnet}
\label{subsection:3.2.1}

Claude Sonnet\footnote{Anthropic, \textit{Claude API Documentation}, \url{https://platform.claude.com/docs/en/home}, truy cập lần cuối ngày 29/05/2026} được sử dụng để xử lý ngôn ngữ tự nhiên và sinh câu trả lời tư vấn. Anthropic cung cấp Claude API và tài liệu tích hợp cho các ứng dụng cần gọi mô hình Claude từ hệ thống bên ngoài. Tài liệu của Anthropic cũng trình bày Claude\footnote{Anthropic, \textit{Claude Models Overview}, \url{https://platform.claude.com/docs/en/about-claude/models/overview}, truy cập lần cuối ngày 29/05/2026} như một họ mô hình ngôn ngữ lớn với nhiều lựa chọn mô hình khác nhau.

Trong hệ thống, AI Service chuẩn bị prompt, dữ liệu truy xuất và ngữ cảnh hội thoại trước khi gọi Claude Sonnet. Cách tổ chức này giúp hệ thống kiểm soát dữ liệu đưa vào mô hình và đảm bảo câu trả lời bám theo knowledge base.

Claude Sonnet được dùng trong các tác vụ: tạo câu trả lời tư vấn, diễn giải kết quả so sánh sản phẩm, giải thích kết quả tham chiếu giá và xử lý hội thoại nhiều lượt.

\subsection{Spring Boot}
\label{subsection:3.2.2}

Spring Boot\footnote{Spring, \textit{Spring Boot Documentation}, \url{https://docs.spring.io/spring-boot/index.html}, truy cập lần cuối ngày 29/05/2026} được sử dụng để xây dựng backend nghiệp vụ. Tài liệu chính thức của Spring Boot mô tả framework này giúp tạo các ứng dụng Spring độc lập, có thể chạy trực tiếp và giảm cấu hình thủ công.

Trong hệ thống, Spring Boot xử lý các nghiệp vụ chính như đăng nhập, phân quyền, quản lý tenant, quản lý thành viên cửa hàng, quản lý chatbot, quản lý nguồn dữ liệu knowledge base, lưu hội thoại, quản lý lead và quản lý yêu cầu mua hàng. Spring Boot cũng là lớp kiểm soát tenant trước khi dữ liệu được truy vấn hoặc cập nhật.

\subsection{FastAPI}
\label{subsection:3.2.3}

FastAPI\footnote{FastAPI, \textit{FastAPI Documentation}, \url{https://fastapi.tiangolo.com/}, truy cập lần cuối ngày 29/05/2026} được sử dụng để xây dựng AI Service. Tài liệu chính thức mô tả FastAPI là framework hiện đại, hiệu năng cao, dùng Python type hints để xây dựng API.

Trong đồ án, FastAPI nhận yêu cầu từ backend, thực hiện truy xuất dữ liệu liên quan, chuẩn bị prompt và gọi Claude Sonnet. Việc tách AI Service khỏi backend nghiệp vụ giúp hệ thống dễ phát triển, kiểm thử và thay đổi phần xử lý AI hơn.

\subsection{PostgreSQL}
\label{subsection:3.2.4}

PostgreSQL\footnote{PostgreSQL Global Development Group, \textit{PostgreSQL Documentation}, \url{https://www.postgresql.org/docs/}, truy cập lần cuối ngày 29/05/2026} được sử dụng làm hệ quản trị cơ sở dữ liệu chính. PostgreSQL là hệ quản trị cơ sở dữ liệu quan hệ mã nguồn mở, có tài liệu chính thức đầy đủ và được sử dụng rộng rãi cho các hệ thống ứng dụng web.

Hệ thống sử dụng PostgreSQL để lưu các dữ liệu như tenant, người dùng, chatbot, nguồn dữ liệu knowledge base, hội thoại, tin nhắn, lead, yêu cầu mua hàng và sản phẩm. Các bảng nghiệp vụ quan trọng được gắn với \texttt{tenant\_id} để phục vụ yêu cầu tách biệt dữ liệu giữa các cửa hàng.

\subsection{HTML, CSS và JavaScript thuần}
\label{subsection:3.2.5}

Frontend của hệ thống được xây dựng bằng HTML\footnote{MDN Web Docs, \textit{HTML: HyperText Markup Language}, \url{https://developer.mozilla.org/en-US/docs/Web/HTML}, truy cập lần cuối ngày 17/06/2026}, CSS\footnote{MDN Web Docs, \textit{CSS: Cascading Style Sheets}, \url{https://developer.mozilla.org/en-US/docs/Web/CSS}, truy cập lần cuối ngày 17/06/2026} và JavaScript\footnote{MDN Web Docs, \textit{JavaScript}, \url{https://developer.mozilla.org/en-US/docs/Web/JavaScript}, truy cập lần cuối ngày 17/06/2026} thuần. HTML được dùng để cấu trúc nội dung trang web. CSS được dùng để trình bày giao diện. JavaScript được dùng để xử lý tương tác phía trình duyệt.

Trong hệ thống, HTML, CSS và JavaScript được dùng cho giao diện quản trị, giao diện tenant và giao diện chat. Giao diện quản trị gồm các màn hình quản lý tenant, thành viên cửa hàng, chatbot, knowledge base, hội thoại, lead và yêu cầu mua hàng. Giao diện chat phục vụ người dùng cuối trong quá trình tư vấn sản phẩm.

\subsection{JWT}
\label{subsection:3.2.6}

Hệ thống có nhiều nhóm người dùng như quản trị hệ thống, quản trị cửa hàng, nhân viên cửa hàng và người dùng cuối. Vì vậy, hệ thống cần cơ chế xác thực và phân quyền rõ ràng. Đồ án sử dụng hướng xác thực dựa trên token, trong đó JWT\footnote{IETF, \textit{RFC 7519: JSON Web Token (JWT)}, \url{https://datatracker.ietf.org/doc/html/rfc7519}, truy cập lần cuối ngày 29/05/2026} là một chuẩn phổ biến để biểu diễn các claim giữa hai bên dưới dạng gọn nhẹ.

Sau khi đăng nhập, người dùng nhận token cho phiên làm việc. Frontend gửi token trong các request tiếp theo để backend kiểm tra danh tính, vai trò và tenant hiện tại. Với các API nghiệp vụ, backend kết hợp kiểm tra token và \texttt{tenant\_id} để đảm bảo người dùng chỉ thao tác trong phạm vi được phép.

\subsection{Webhook}
\label{subsection:3.2.7}

Webhook được sử dụng để tích hợp các kênh chat ngoài như Messenger\footnote{Meta for Developers, \textit{Graph API Webhooks}, \url{https://developers.facebook.com/docs/graph-api/webhooks/}, truy cập lần cuối ngày 29/05/2026} hoặc Telegram\footnote{Telegram, \textit{Telegram Bot API}, \url{https://core.telegram.org/bots/api}, truy cập lần cuối ngày 29/05/2026}.

Trong hệ thống, webhook tiếp nhận tin nhắn từ kênh ngoài, xác định chatbot hoặc tenant tương ứng, sau đó chuyển nội dung vào luồng xử lý hội thoại. Câu trả lời được tạo thông qua AI Service và gửi lại cho người dùng qua API của nền tảng chat.

\subsection{Docker và Docker Compose}
\label{subsection:3.2.8}

Hệ thống gồm nhiều thành phần như frontend, backend, AI Service và PostgreSQL. Docker và Docker Compose\footnote{Docker, \textit{Docker Compose Documentation}, \url{https://docs.docker.com/compose/}, truy cập lần cuối ngày 29/05/2026} được sử dụng để đóng gói và khởi chạy các thành phần này trong môi trường demo. Docker Compose cho phép định nghĩa và chạy ứng dụng gồm nhiều container thông qua một tệp cấu hình.

\section{Tổng hợp lựa chọn công nghệ}
\label{section:3.3}

Bảng~\ref{table:technology-summary} tổng hợp ngắn gọn vai trò của các công nghệ chính được sử dụng trong đồ án.

\begin{table}[H]
\centering
\caption{Tổng hợp công nghệ sử dụng trong hệ thống}
\label{table:technology-summary}
\resizebox{\textwidth}{!}{
\begin{tabular}{|p{3.5cm}|p{4.2cm}|p{7cm}|}
\hline
\textbf{Thành phần} & \textbf{Công nghệ sử dụng} & \textbf{Vai trò trong hệ thống} \\
\hline
Backend nghiệp vụ & Spring Boot & Xử lý API, phân quyền, tenant, chatbot, knowledge base, hội thoại, lead và yêu cầu mua hàng. \\
\hline
AI Service & FastAPI, Claude Sonnet & Truy xuất dữ liệu, chuẩn bị prompt, sinh câu trả lời tư vấn, so sánh và tham chiếu giá. \\
\hline
Cơ sở dữ liệu & PostgreSQL & Lưu dữ liệu tenant, người dùng, sản phẩm, hội thoại, lead và yêu cầu mua hàng. \\
\hline
Frontend & HTML, CSS và JavaScript thuần & Xây dựng giao diện quản trị và giao diện chat dưới dạng static resources được Spring Boot phục vụ trực tiếp. \\
\hline
Xác thực/phân quyền & JWT & Xác thực người dùng, truyền thông tin vai trò và tenant trong request. \\
\hline
Tích hợp kênh ngoài & Webhook & Nhận tin nhắn từ Messenger/Telegram và đưa vào luồng xử lý chatbot. \\
\hline
Triển khai demo & Docker Compose & Khởi chạy các service của hệ thống trong môi trường thống nhất. \\
\hline
\end{tabular}
}
\end{table}

\section{Kết chương}
\label{section:3.ket_chuong}

Chương này đã trình bày các cơ sở lý thuyết và công nghệ chính được sử dụng trong đồ án. Về lý thuyết, hệ thống dựa trên mô hình multi-tenant để phục vụ nhiều cửa hàng, sử dụng RAG và knowledge base để chatbot tư vấn theo dữ liệu riêng, đồng thời tổ chức dữ liệu sản phẩm để hỗ trợ so sánh và tham chiếu giá.

Về công nghệ, hệ thống sử dụng Spring Boot cho backend nghiệp vụ, FastAPI cho AI Service, Claude Sonnet cho sinh câu trả lời, PostgreSQL cho lưu trữ dữ liệu, HTML/CSS/JavaScript thuần cho frontend, JWT cho xác thực/phân quyền, webhook cho tích hợp kênh chat ngoài và Docker Compose cho môi trường demo.

Các nội dung trong chương này là cơ sở để trình bày thiết kế kiến trúc, mô hình dữ liệu và luồng xử lý hệ thống ở các chương tiếp theo.

\end{document}